%-------------------------
% Resume in Latex
% Author : Jake Gutierrez [Customized by Raghav Senthil Kumar]
% License : MIT
%------------------------

\documentclass[letterpaper,10pt]{article}

% Shared LaTeX preamble for every resume (packages, margins, section style, custom commands).
% Compile from the .tex file's own directory, e.g.  cd base && pdflatex <file>.tex
% Edit formatting HERE ONCE; do not copy it back into individual resumes.

\usepackage{latexsym}
\usepackage[empty]{fullpage}
\usepackage{titlesec}
\usepackage{marvosym}
\usepackage[usenames,dvipsnames]{color}
\usepackage{verbatim}
\usepackage{enumitem}
\usepackage[hidelinks]{hyperref}
\usepackage{fancyhdr}
\usepackage[english]{babel}
\usepackage{tabularx}
\IfFileExists{glyphtounicode.tex}{\input{glyphtounicode}}{}



\pagestyle{fancy}
\fancyhf{}
\fancyfoot{}
\renewcommand{\headrulewidth}{0pt}
\renewcommand{\footrulewidth}{0pt}

% Adjust margins
\addtolength{\oddsidemargin}{-0.5in}
\addtolength{\evensidemargin}{-0.5in}
\addtolength{\textwidth}{1in}
\addtolength{\topmargin}{-.5in}
\addtolength{\textheight}{1.0in}
\setlength{\footskip}{4.08003pt}

\urlstyle{same}

\raggedbottom
\raggedright
\setlength{\tabcolsep}{0in}

% Sections formatting
\titleformat{\section}{
  \vspace{-4pt}\scshape\raggedright\large
}{}{0em}{}[\color{black}\titlerule \vspace{-5pt}]

% Ensure that generate pdf is machine readable/ATS parsable
\ifdefined\pdfgentounicode
  \pdfgentounicode=1
\fi

%-------------------------
% Custom commands
\newcommand{\resumeItem}[1]{
  \item{
    {#1}
  }
}

\newcommand{\resumeSubheading}[4]{
  \vspace{-2pt}\item
    \begin{tabular*}{0.97\textwidth}[t]{l@{\extracolsep{\fill}}r}
      \textbf{#1} & #2 \\
      \textit{#3} & \textit{#4} \\
    \end{tabular*}\vspace{-7pt}
}

\newcommand{\resumeSubSubheading}[2]{
    \item
    \begin{tabular*}{0.97\textwidth}{l@{\extracolsep{\fill}}r}
      \textit{\small#1} & \textit{\small #2} \\
    \end{tabular*}\vspace{-7pt}
}

\newcommand{\resumeProjectHeading}[2]{
    \item
    \begin{tabular*}{0.97\textwidth}{l@{\extracolsep{\fill}}r}
      #1 & #2 \\
    \end{tabular*}\vspace{-7pt}
}

\newcommand{\resumeSubItem}[1]{\resumeItem{#1}\vspace{-4pt}}

\renewcommand\labelitemii{$\vcenter{\hbox{\tiny$\bullet$}}$}

\newcommand{\resumeSubHeadingListStart}{\begin{itemize}[leftmargin=0.15in, label={}]}
\newcommand{\resumeSubHeadingListEnd}{\end{itemize}}
\newcommand{\resumeItemListStart}{\begin{itemize}[leftmargin=0.22in, itemsep=1pt, topsep=2pt, parsep=0pt, partopsep=0pt]}
\newcommand{\resumeItemListEnd}{\end{itemize}\vspace{-2pt}}


\begin{document}

\begin{center}
    {\Huge\bfseries Raghav Senthil Kumar} \\ \vspace{1pt}
    \small 623-332-2168 $|$ \href{mailto:raghavs9@illinois.edu}{\underline{raghavs9@illinois.edu}} $|$
    \href{https://linkedin.com/in/raghav-senthil-kumar}{\underline{linkedin.com/in/raghav}} $|$
    \href{https://github.com/Raghavsk24}{\underline{github.com/raghav}}
\end{center}


%-----------EDUCATION-----------
\section{Education}
  \resumeSubHeadingListStart
    \resumeSubheading
      {University of Illinois Urbana-Champaign}{Champaign, IL}
      {B.S. in Computer Science and Statistics}{Expected: May 2029}
      \resumeItemListStart
        \resumeItem{\textbf{Coursework:} Computer Science I (Java), Computer Science II (C++), Data Science \& ML (Python), Calculus III}
        \resumeItem{\textbf{Activities:} Disruption Lab, SIGRobotics, AI Alignment, BuildIllinois, Innovation LLC, Agentic AI}
      \resumeItemListEnd
  \resumeSubHeadingListEnd


%-----------EXPERIENCE-----------
\section{Experience}
  \resumeSubHeadingListStart

    \resumeProjectHeading
      {\textbf{Founding Software Engineer $|$ Clusion}}{September 2025 -- January 2026}
      \resumeItemListStart
        \resumeItem{Built a SaaS platform to automate research \& case-writing in competitive high school debate, reaching \textbf{10K+ users}}
        \resumeItem{Designed a text retrieval algorithm that extracts evidence segments from a debate case using URL validation \& IQR filtering. Indexed \textbf{600K+ documents} on an \emph{Elasticsearch} engine, returning query results in \textbf{less than 1 second}}
        \resumeItem{Developed a de-duplication algorithm that flagged \textbf{30\%} of evidence documents as identical, reducing index bloat}
      \resumeItemListEnd

    \resumeProjectHeading
      {\textbf{Product Development Intern $|$ Jetson}}{January 2025 -- February 2025}
      \resumeItemListStart
        \resumeItem{Conducted \textbf{7 iterations of UI/UX testing} across 3 experimental features on the Jetson App (100K+ users)}
        \resumeItem{Delivered a \textbf{5-page report} on recommending strategies to optimize user engagement directly to the CEO}
        \resumeItem{Redesigned user onboarding on \emph{Figma}, cutting it from \textbf{12 to 5 pages}, validated via \textbf{A/B testing} on 30 test users}
      \resumeItemListEnd

    \resumeProjectHeading
      {\textbf{Data Analytics Intern $|$ STEM E}}{June 2024 -- July 2024}
      \resumeItemListStart
        \resumeItem{Analyzed 3 years of TikTok data (\textbf{200+ posts}) on \emph{Hootsuite} and \emph{Microsoft Excel} to find patterns in outreach}
        \resumeItem{Identified that ``science fun facts'' posts comprised 69.3\% of views, despite accounting for 7\% of all content}
        \resumeItem{Pitched new content strategy centered around ``fun facts'' that \textbf{increased engagement by 24\% (8K+ views)}}
      \resumeItemListEnd

  \resumeSubHeadingListEnd


%-----------PROJECTS-----------
\section{Projects}
    \resumeSubHeadingListStart
    \resumeProjectHeading
          {\textbf{Truck Pricing VLM} $|$ \href{https://truck-pricing-vlm.vercel.app/}{\underline{Live Website}} $|$ \href{https://github.com/Raghavsk24/Truck-Pricing-VLM}{\underline{Source Code}}}{September 2026}
          \resumeItemListStart
            \resumeItem{Led team of 4 to build a VLM-based pricing algorithm from user-inputted truck images for \textbf{Kamion (YC S22)}}
            \resumeItem{Constructed dataset of \emph{15,000 images} scraped from publicly-available truck listings and audited for data quality}
            \resumeItem{Fine-tuned \emph{Claude Sonnet 4.6 VLM} across \textbf{30 independent iterations} \& wrote JSON schema for feature extration}
            \resumeItem{Optimized pricing algorithm via \emph{quantile regression} with \textbf{87\% accuracy}. \textbf{\underline{2nd at UIUC Founder's Hackathon}}}
          \resumeItemListEnd
      \resumeProjectHeading
          {\textbf{SchoolPulse AI} $|$ \href{https://school-pulse-ai.vercel.app/}{\underline{Live Website}} $|$ \href{https://github.com/vignesh-nagarajan-vn/SchoolPulse-AI}{\underline{Source Code}}}{June 2026}
          \resumeItemListStart
            \resumeItem{Engineered a \textbf{suite of 3 edge AI tools} that collects data regarding a school's energy, water and food consumption}
            \resumeItem{Built a \emph{Python FastAPI} backend that ingests live water, waste \& energy data from all 3 tools, runs \textbf{4 Random Forest models} to detect anomalies \& deliver actionable steps via \textbf{\emph{Gemma} Voice LLM} on a \emph{Next.js} dashboard}
            \resumeItem{Led a team of 5 across project ideation, development and presentation. \textbf{\underline{Finalist at USAII Global Hackathon}}}
          \resumeItemListEnd
      \resumeProjectHeading
          {\textbf{EcoBin} $|$ \href{https://arxiv.org/abs/2606.15547}{\underline{Research Paper}} $|$ \href{https://github.com/Raghavsk24/EcoBin}{\underline{Source Code}}}{March 2026 -- May 2026}
          \resumeItemListStart
            \resumeItem{Designed a two-stage \emph{EfficientNetV2-S} neural network (\textbf{4.39M} parameters) trained on \textbf{15K+ images} that classifies waste into 4 disposal pathways, achieving \textbf{96.13\% overall accuracy}. Published 2 original datasets onto \emph{Kaggle}}
            \resumeItem{Built a synthetic data generation pipeline using \emph{U2-Net} segmentation and alpha-compositing to generate 9K images to train the model to detect contamination in recyclable waste. \textbf{\underline{Placed 2nd at the Arizona State Science Fair}}}
            \resumeItem{Engineered a smart bin prototype using \emph{Arduino Uno} that sorts the waste item with \textbf{less than 2 second latency}}
          \resumeItemListEnd
      \resumeProjectHeading
          {\textbf{Chronos} $|$ \href{https://usechronos.live/}{\underline{Live Website}} $|$ \href{https://github.com/Raghavsk24/Chronos}{\underline{Source Code}}}{December 2025}
          \resumeItemListStart
            \resumeItem{Built a \textbf{scheduling agent} (\emph{React/Firebase}) that auto-books meetings for large groups (30+ ppl) on Google Calendar}
            \resumeItem{Designed a \textbf{5-step \emph{Python} scheduling algorithm} that scans calendar availability, returns a ranked list of time slots and sends email confirmation for each booking. Validated algorithm across \textbf{14 unit and integration test cases}}
            \resumeItem{Created a lobby system that enables the host to authenticate invited participants and queue follow-up meetings}
            \resumeItem{Licensed Chronos to \textbf{8 startups} \& signed \textbf{2 LOIs worth \$3.5K}. \textbf{\underline{Won 2025 Congressional App Challenge}}}
          \resumeItemListEnd
    \resumeSubHeadingListEnd



%
%-----------TECHNICAL SKILLS-----------
\section{Technical Skills}
 \begin{itemize}[leftmargin=0.15in, label={}]
    \item{
     \textbf{Programming Languages}{: Java, Python, C++, JavaScript, SQL, HTML/CSS} \\
     \textbf{Frameworks/Platforms}{: FastAPI, React, Next.js, Node.js, PostgreSQL, Elasticsearch, Supabase, Firebase, AWS} \\
     \textbf{DevTools}{: Git/GitHub, Docker, CI/CD, Vercel, Google Cloud, VS Code, Claude Code, Codex, Figma} \\
     \textbf{AI/ML}{: PyTorch, TensorFlow, Keras, scikit-learn, OpenCV, Hugging Face, pandas, NumPy, Claude API} \\
     \textbf{Concepts}{: Object-Oriented Programming, REST APIs, Unit Testing, ETL, LLMs, RAG, AI Agents, Agile}
    }
 \end{itemize}
 \vspace{3.53pt}

%-------------------------------------------
\end{document}
